\section{El modelo como nuevo sistema operativo}

La entrevista plantea que «el modelo es el sistema operativo de nueva generación». Esto no significa que el modelo sustituya a Linux, Windows o iOS, sino que el modelo podría convertirse en una nueva capa de orquestación entre la intención del usuario y la ejecución de las aplicaciones. Antes, el OS administraba el hardware, los archivos, los procesos y las aplicaciones; en el futuro, el modelo/Agent OS podría administrar intenciones, herramientas, contexto, permisos y el estado de las tareas.

\begin{figure}[H]
\centering
\includegraphics[width=0.92\textwidth]{model-as-os.png}
\caption{El modelo como OS de nueva generación: comprensión de la intención, orquestación de herramientas y memoria del estado.}
\end{figure}

\begin{knowledgebox}{Lectura de la figura: qué administra el Model OS}
En la figura, el OS del modelo ocupa el centro y conecta la intención del usuario, los servicios de las aplicaciones, la memoria de datos, el entorno de ejecución y los permisos de seguridad. La tesis que respalda es esta: el modelo no es solo una app, sino que podría convertirse en el nuevo punto de entrada a las aplicaciones y en la capa de orquestación de tareas.
\end{knowledgebox}

\subsection{Por qué Coding es la puerta de entrada para que el modelo funcione como OS}

El código es el lenguaje formal más potente del mundo digital. Si el modelo puede escribir código, modificarlo, llamar a API y ejecutar scripts, puede rebasar los límites de una sola app y convertirse en una capa intermedia que opera varios sistemas. Por eso, Coding es la puerta de entrada decisiva para que el modelo funcione como OS.

\begin{importantbox}{El cambio de app a OS}
Una app es un punto de entrada a una sola función; un OS es una capa de orquestación de recursos y tareas. Si el modelo puede comprender intenciones, llamar a herramientas, administrar el contexto y ejecutar tareas, pasará de la posición de app a la de OS.
\end{importantbox}

\subsection{Los tres grandes de China y el OS del modelo}

La entrevista menciona a los tres grandes de China, pero lo importante no es la clasificación concreta, sino una tesis estructural: si los equipos de modelos del país quieren aprovechar el segundo acto, necesitan comprender a la vez Coding, Agent, costos y escenarios de aplicación. La competencia por el OS del modelo no consiste solo en hacer un producto de chat, sino en integrarse con el trabajo de oficina, el desarrollo, el contenido, los datos, la automatización y los procesos empresariales.

\subsection{Resumen de la sección}

Que el modelo se convierta en el OS de nueva generación significa que pasa de ser un sistema que responde a ser un sistema que ejecuta tareas. Coding y el harness son la vía decisiva para que llegue a esa posición.
